\subsection{Decision value with binding task constraints}
\label{sec:r16activefacevalue}
Temporary action constraints need not rule out a strict decision-value result. The interior assumption in Proposition~\ref{prop:r16decisionvalue} can be replaced by a stable-face condition. The relevant prediction norm then charges errors only in the directions in which each economic task can respond.

Retain the quadratic objectives and the common observation $z=q+\varepsilon$ of Section~\ref{sec:r16decisionvalue}, with $\E\varepsilon=0$ and covariance $\Sigma$. For each task, fix a nonempty face $F_j$ of its feasible polytope, a point $a_j^0\in F_j$, and a full-column-rank matrix $L_j$ spanning the parallel space of $F_j$. Put
\[
 R_j=L_j(L_j^\top H_jL_j)^{-1}L_j^\top,
 \qquad W_F=\sum_{j=1}^K\omega_jh_j^2R_j.
\]
For a singleton face set $R_j=0$. Assume that $W_F$ is positive definite. Individual faces may have binding constraints; the last condition requires their decision-relevant tangent spaces jointly to span the continuation coordinates.

Let $D$ be a fixed scalar-gradient dictionary, independent of $\varepsilon$, and let $P_F$ project orthogonally onto the columns of $W_F^{1/2}D$. Define
\[
 \widehat q_R=z,\qquad
 \widehat q_C=W_F^{-1/2}P_FW_F^{1/2}z,
 \qquad y_F=W_F^{1/2}q,
 \quad\Omega_F=W_F^{1/2}\Sigma W_F^{1/2}.
\]
The stable-face condition is that, for $q$, $\widehat q_R$, and $\widehat q_C$, the unique constrained maximizer lies in the relative interior of the same deterministic face $F_j$, almost surely, for every task. This is a condition on the original observation law; it does not condition an otherwise unrestricted noise distribution on a face-selection event. It can be verified by the inactive slacks and the normal optimality inequalities on the prescribed noise support.

\begin{proposition}[Continuation compression on stable active faces]
\label{prop:r16activefacevalue}
Under these conditions, the expected weighted payoff difference between the constrained critic and Raw actions equals
\begin{equation}
 \frac12\left[\operatorname{tr}\{(I-P_F)\Omega_F\}
                  -\|(I-P_F)y_F\|^2\right].
 \label{eq:r16activefacevalue}
\end{equation}
Consequently, compression strictly improves expected economic payoff whenever the removed decision-relevant variance exceeds the weighted approximation cost. In particular, this holds for a correctly specified dictionary that removes strictly positive decision-relevant variance, even though every task may have binding constraints.
\end{proposition}
\begin{proof}
Tangential optimality on $F_j$ gives, for each of the three continuation vectors $v$,
\[
 a_j(v)=a_j^0+R_j(b_j-H_ja_j^0-h_jv).
\]
Thus $a_j(v)-a_j(q)=-h_jR_j(v-q)$. The gradient of the true objective at $a_j(q)$ is orthogonal to this difference. Expanding the quadratic objective, and using $R_jH_jR_j=R_j$, yields
\[
 Q_j(a_j(q);q)-Q_j(a_j(v);q)
       =\tfrac12 h_j^2(v-q)^\top R_j(v-q).
\]
Summing the weighted losses gives Raw expected loss
\[
 \tfrac12\operatorname{tr}(\Omega_F).
\]
For the critic, the weighted error is
\[
 -(I-P_F)y_F+P_FW_F^{1/2}\varepsilon.
\]
Its expected squared norm is
\[
 \|(I-P_F)y_F\|^2+\operatorname{tr}(P_F\Omega_F).
\]
Subtracting the two expected losses proves the formula and its strictness assertions.
\end{proof}

The result gives a positive boundary-aware mechanism, not a universal ranking of methods. The faces and dictionary must satisfy the stated conditions for the full observation law. If fitted actions switch faces, the displayed exact identity cannot be obtained merely by selecting the observations with a favorable face. If $W_F$ is singular, the present statement does not invoke an ordinary inverse; a lower-dimensional formulation requires a separately specified decision space. The nonlinear menu is evaluated by its own paired payoff intervals, rather than by asserting that its fitted features or active sets satisfy this specialization. Identical projected estimators, whether called Raw or NBO, still have identical decisions.
